\documentclass{article}
\usepackage{mathtools}
\usepackage{graphicx}
\usepackage{xcolor}
\usepackage{url}
\usepackage[margin=1in]{geometry}
\usepackage{filecontents}


\begin{document}
\title{Peer Review - Julia \\
Assessment of Artificial Intelligence in Echocardiography Diagnostics in Differentiating Takotsubo Syndrome From Myocardial Infarction}
\author{Dustin Randall}
\date{}
\maketitle

\section{What was interesting}
I found this article interesting because it was able to explore the idea of applying a machine learning model to 2d gray scale echocardiograms.
By leveraging a multi-step approach to process the data, they were able to improve the accuracy and reliability of the model's ability to distinguish between the two symptoms.
I also found the self-improving model approach full of potential for many different applications.

\section{Why is it important}
This research is important because it provides a mechanism to diagnose more patients efficiently and more cheaply.
This can make optimal use of limited medical resources by focusing their time and efforts toward the right mitigations.
Echocardiograms are much less invasive and more accessible, and training an AI model to interpret the results can significantly improve diagnostic throughput.

\section{Presentation}
I thought that the presentation was very well done.
The visuals were engaging, and I felt like I was able to follow along with the background and purpose of the paper.
Using machine learning to improve the state of the art in medicine seems inherently important, and the presentation helped make it interesting as well.

\section{Citations}
Contemporary, non-invasive imaging diagnosis of chronic coronary artery disease~\cite{bijl_contemporary_2025}
and {2D} echocardiography video to {3D} heart shape reconstruction for clinical application~\cite{laumer_2d_2025}

\bibliography{citations}
\bibliographystyle{plain}

\end{document}